\section{Parameterized enumeration and asymptotic regimes}
\label{sec:enumeration}

\subsection{The birth-front search}

The normal form leads to a direct iterative-deepening algorithm.

\begin{algorithmic}[Birth-front unlocking search]
\label{alg:birth-front}
Given a state $x$, a depth bound $k$, and a birth bound $b$:
\begin{enumerate}
\item For $\ell=1,\ldots,k$, search traces of length at most $\ell$.
\item Enumerate a canonical list of at most $b$ pairwise footprint-disjoint
      births.  Each birth is chosen among all currently legal neutral actions.
\item After the birth block, permit only a legal action whose footprint meets
      the union of all previously chosen footprints.
\item After every neutral action, test for a terminal reduction.  Return the
      shortest successful trace.
\item Undo failed branches exactly.  Exhaustion returns \emph{unknown}, never a
      negative verdict.
\end{enumerate}
\end{algorithmic}

A canonical order on births removes permutations of pairwise commuting actions.
The complexity theorem below does not need that saving, so its bound is robust
to a simpler implementation.

\subsection{Counting search nodes}

\begin{theorem}[Parameterized enumeration]
\label{thm:enumeration}
Assume footprint size at most $f$, global incidence at most $AN$, local
incidence at most $B|S|$, and polynomial-time state operations.  The birth-front
algorithm finds a first-unlocking trace whenever one of length at most $k$ and
birth number at most $b$ exists, and exhausts all birth-front normal-form
candidates in time
\[
   N^{b+O(1)}(Ck)^k,
   \qquad C=Bf\max(1,A),
\]
and polynomial working memory.  The constant implicit in $O(1)$ is independent
of $b$ and $k$.
\end{theorem}

\begin{proof}
By Theorem~\ref{thm:birth-front}, it is enough to enumerate traces with all
births first.  Suppose a normal-form trace has $r\leq b$ births and total length
$\ell\leq k$.

At each birth position there are at most $AN$ legal neutral actions, so the birth
block has at most $(AN)^r$ possibilities.  After $j$ actions have been chosen,
the accumulated footprint contains at most $fj$ sites.  Bounded local incidence
therefore gives at most $Bfj$ possible connected extensions.  Consequently,
the number of length-$\ell$ search words with $r$ births is at most
\[
  (AN)^r\prod_{j=r}^{\ell-1} Bfj
  \leq (AN)^b(Bfk)^k.
\]
Summing over $r\leq b$ and $\ell\leq k$ contributes at most a polynomial or
$2^{O(k)}$ factor, which is absorbed by replacing $C$ with a larger absolute
rule-set constant.

Each node performs polynomial-time legality, application, terminal, and undo
operations.  A depth-first traversal stores one state representation, at most
$k$ undo records, the accumulated footprint of size at most $fk$, and polynomial
bookkeeping.  Hence the working memory is polynomial in $N$ and $k$.
\end{proof}

\begin{remark}[Why the exponent $b$ appears]
Every birth can be located anywhere in the diagram, so a direct deterministic
search pays one factor of $N$ per independent branch.  Removing this exponent
would require extra structure, such as color-coding, representative families,
a reverse terminal-rooted search, or a geometric theorem constraining where
births can occur.
\end{remark}

\subsection{Quasi-polynomial and polynomial corollaries}

Write
\[
 T(N;b,k)=N^{b+O(1)}(Ck)^k.
\]
Taking logarithms gives
\[
 \log T=(b+O(1))\log N + k\log(Ck).
\]

\begin{corollary}[Logarithmic parameters]
\label{cor:log-parameters}
If $b=O(\log N)$ and $k=O(\log N)$, then
\[
 T(N;b,k)=\exp(O((\log N)^2))=N^{O(\log N)}.
\]
Thus the search is quasi-polynomial.
\end{corollary}

\begin{corollary}[Constant branching]
\label{cor:constant-branching}
If $b=O(1)$ and $k=O(\log N)$, then
\[
 T(N;b,k)=\exp(O(\log N\log\log N))
          =N^{O(\log\log N)}.
\]
\end{corollary}

\begin{corollary}[Sublogarithmic depth]
\label{cor:polynomial-island}
If $b=O(1)$ and
$k=O(\log N/\log\log N)$, then $T(N;b,k)=N^{O(1)}$.
\end{corollary}

\begin{proof}
For the last statement,
$k\log(Ck)=O(\log N)$, while the birth factor is polynomial because $b$ is
constant.
\end{proof}

\subsection{Search budgets}

A practical implementation will cap the number of trial moves rather than run
the worst-case tree.  The theorem and the budget serve different purposes.

\begin{itemize}
\item The theorem says what is complete when the specified $(b,k)$ tree is
      exhausted.
\item A trial budget may stop earlier to protect the exact fallback.
\item Budget exhaustion has no mathematical meaning about knottedness or about
      the nonexistence of a trace.
\item Successful traces remain exact because they are independently replayed.
\end{itemize}

This is compatible with the current ProveIt design, which already treats RIII
search as an optional pre-scan accelerator.
